\hspace{3mm}

\centerline{\bf \huge 6 класс}

\begin{flushright}
    \scriptsize{ $-$ Если сдаваться, не попробовав, результат не изменится.}
\end{flushright}

\problem{1} Среди 20 учеников ЭлМШ есть шиншиллы, которые всегда говорят правду, и лососи, которые всегда лгут. Какое может быть наибольшее количество лососей, если каждый сказал, что всего среди 20 учеников не более 2 шиншилл.

\hspace{3mm}

\problem{2} Однажды Очир решил отправиться из Элисты в разностороннее путешествие. На своё усмотрение за одно действие он может поехать либо строго на юг, либо строго на север. За первое действие Очир проехал 1км, а за каждое последующее действие в два раза больше предыдущего(2км, 4км, 8км, 16км и так далее). Может ли Очир вернуться в Элисту за некоторое количество таких действий (Элиста не объемная и является точкой на плоскости).


\hspace{3mm}


\problem{3} Фёдор пришел на I ежегодную Математическую Олимпиаду ЭлМШ, на которой было предложено 5 задач. Фёдор очень торопился и мог потратить на Олимпиаду максимум час. Задачи, придуманные Расулом Владимировичем, Фёдор читает за 2 минуты, а решает за 35 минут. Задачи, придуманные Анджуром Аркадьевичем, он читает 2 минуты, а решает за 15 минут. Задачи, придуманные Очиром Владимировичем, он читает 10 минут и решает тоже за 10 минут. Очир Владимирович предложил две задачи, Анджур Аркадьевич одну, а Расул Владимирович две. Какое наибольшее количество баллов он успел набрать, если задачи Очира Владимировича он мог отличить, даже не читая, а отличить задачи Анджура Аркадьевича от задач Расула Владимировича может только прочитав их, при этом, читая, сначала ему попадаются задачи только Расула Владимировича. За задачи Очира Владимировича ученики получали по 4 балла, Анджура Аркадьевича - по 5 баллов, Расула Владимировича - по 8 баллов.

\hspace{3mm}


\problem{4} После олимпиады Герман решил подшутить над Расулом Владимировичем и Анджуром Аркадьевичем и подговорил всех учеников ЭлМШ прийти на апелляцию. Шиншиллы по своей неопытности не умели лить воду, поэтому просто молчали на апелляции. Лососи из-за недостаточного количества опыта лили воду любому проверяющему. Киви же, зная учителей, лили воду только Анджуру Аркадьевичу. 

\newpage

\hspace{3mm}

\noindentОказалось, что олимпиаду проверили просто идеально, проставив всем только заслуженные баллы. Но Анджур Аркадьевич все равно добавлял по баллу тем, кто уверенно лил воду. Расул Владимирович же, наоборот, отнимал балл у тех, кто пытался лить воду. После апелляции оказалось, что баллов стало на 15 больше. Сколько Шиншилл апеллировалось у Расула Владимировича, если количество попавших к нему Киви было столько же, сколько Шиншилл попало к Анджуру Аркадьевичу. Причем Анджур Аркадьевич и Расул Владимирович приняли одинаковое количество детей.

\hspace{3mm}


\problem{5} Божья коровка и муравей устроили забег на часах. Муравей бежал по минутной стрелке часов, а божья коровка - по часовой стрелке. Когда впервые с начала гонки стрелки на часах пересеклись, муравей пробежал $\frac{2}{5}$ дистанции, а божья коровка только $\frac{1}{4}$. Когда божья коровка завершила забег, если муравей добежал до конца стрелки в 12:00? Забег начался одновременно, а божья коровка и муравей бежали с постоянной скоростью.